% Zone „Das kennst du schon“ – aus bank/ableitungsregeln/zone.jsonl (zusammenbau v0.1)
\einheitenkopf[zone][Kennst du schon]{Das kennst du schon}
\setcounter{aufgabe}{0}

%% TODO Titel: Ich-kann-Satz fehlt in der Bank (Platzhalter: Kettenname) – Nr. 1
%% TODO Anweisung: ein Satz über der ersten Teilaufgabe fehlt in der Bank – Nr. 1
\begin{aufgabe}{Potenzen mit natürlichem Exponenten lesen und schreiben, Vorfaktor und Exponent unterscheiden, Potenzgesetze (Produkt gleicher Basen, Potenz einer Potenz), negative Exponenten als Kehrwert, Wurzel als Potenz mit Exponent ein Halb}
\begin{teile}
\teil $x^3 \cdot x^5$ als eine Potenz? \leerfeld
\teil $\frac{1}{x^5}$ als Potenz von $x$? \leerfeld
\end{teile}
\end{aufgabe}

%% TODO Titel: Ich-kann-Satz fehlt in der Bank (Platzhalter: Kettenname) – Nr. 2
%% TODO Anweisung: ein Satz über der ersten Teilaufgabe fehlt in der Bank – Nr. 2
\begin{aufgabe}{Terme umformen}
\begin{teile}
\teil $4(x + 5)$ ausmultipliziert? \leerfeld
\teil $5x \cdot e^x - 2x^2 \cdot e^x$ – $e^x$ ausgeklammert? \leerfeld
\end{teile}
\end{aufgabe}

%% TODO Titel: Ich-kann-Satz fehlt in der Bank (Platzhalter: Kettenname) – Nr. 3
%% TODO Anweisung: ein Satz über der ersten Teilaufgabe fehlt in der Bank – Nr. 3
\begin{aufgabe}{Bruch- und Dezimalkoeffizienten multiplizieren (ein Viertel mal fünf, null Komma vier mal vier), mit Brüchen kürzen}
\begin{teile}
\teil $\frac{1}{3} \cdot 6$? \leerfeld
\teil $\frac{3}{8} \cdot 4$ – gekürzt? \leerfeld
\end{teile}
\end{aufgabe}

%% TODO Titel: Ich-kann-Satz fehlt in der Bank (Platzhalter: Kettenname) – Nr. 4
%% TODO Anweisung: ein Satz über der ersten Teilaufgabe fehlt in der Bank – Nr. 4
\begin{aufgabe}{Die e-Funktion als Funktion, die sich selbst als Ableitung hat; $e^{(kx)}$ als Verkettung lesen; $e^x$ ist nie null und immer positiv}
\begin{teile}
\teil Welchen Wert hat $e^0$? \leerfeld
\teil $e^{-0{,}25x}$ an der Stelle $x = 8$ – Wert auf zwei Stellen gerundet? \leerfeld
\end{teile}
\end{aufgabe}

%% TODO Titel: Ich-kann-Satz fehlt in der Bank (Platzhalter: Kettenname) – Nr. 5
%% TODO Anweisung: ein Satz über der ersten Teilaufgabe fehlt in der Bank – Nr. 5
\begin{aufgabe}{Ableitung als Tangentensteigung und Änderungsrate deuten, Ableitungswert an einer Stelle als Steigung lesen, Extrempunkt als Stelle mit waagerechter Tangente}
\begin{teile}
\teil Lies die Steigung der Tangente $t$ an der Stelle $x = 2$ ab.

\begin{ksys}[xmin=-1,xmax=5,ymin=-1,ymax=5,ablesen] \funktion{0.25*\x^2}{f} \tangentean{0.25*\x^2}{2}{t} \end{ksys}
Steigung = \leerfeld
\teil $f'(1) = -1{,}5$ – steigt oder fällt der Graph von $f$ bei $x = 1$, und wie steil? \leerfeld
\end{teile}
\end{aufgabe}

%% TODO Titel: Ich-kann-Satz fehlt in der Bank (Platzhalter: Kettenname) – Nr. 6
%% TODO Anweisung: ein Satz über der ersten Teilaufgabe fehlt in der Bank – Nr. 6
\begin{aufgabe}{Funktionswerte berechnen und Punktprobe, Anstieg an einer Stelle als Wert der ersten Ableitung}
\begin{teile}
\teil $f(x) = x^2 + 4$ – $f(3)$? f(3) = \leerfeld
\teil $f'(x) = 2x^2 - 5$ – Anstieg des Graphen an der Stelle $x = -2$? f'(-2) = \leerfeld
\end{teile}
\end{aufgabe}

%% TODO Titel: Ich-kann-Satz fehlt in der Bank (Platzhalter: Kettenname) – Nr. 7
%% TODO Anweisung: ein Satz über der ersten Teilaufgabe fehlt in der Bank – Nr. 7
\begin{aufgabe}{Verschiebung und Streckung eines Graphen an der Gleichung erkennen (f(x − c) verschiebt nach rechts, c · f(x) streckt)}
\begin{teile}
\teil Wie entsteht der Graph von $y = (x - 4)^2$ aus dem von $y = x^2$? \leerfeld
\teil Wie entsteht der Graph von $y = (x + 2{,}5)^2$ aus dem von $y = x^2$? \leerfeld
\end{teile}
\end{aufgabe}

%% TODO Titel: Ich-kann-Satz fehlt in der Bank (Platzhalter: Kettenname) – Nr. 8
%% TODO Anweisung: ein Satz über der ersten Teilaufgabe fehlt in der Bank – Nr. 8
\begin{aufgabe}{Potenzen mit natürlichem Exponenten lesen und schreiben, Vorfaktor und Exponent unterscheiden, Potenzgesetze (Produkt gleicher Basen, Potenz einer Potenz), negative Exponenten als Kehrwert, Wurzel als Potenz mit Exponent ein Halb – Fehler finden}
\begin{teile}
\teil Ich finde den Fehler: Tom rechnet $(2x)^4 = 2x^4$.
\end{teile}
\end{aufgabe}

%% TODO Titel: Ich-kann-Satz fehlt in der Bank (Platzhalter: Kettenname) – Nr. 9
%% TODO Anweisung: ein Satz über der ersten Teilaufgabe fehlt in der Bank – Nr. 9
\begin{aufgabe}{Potenzen mit natürlichem Exponenten lesen und schreiben, Vorfaktor und Exponent unterscheiden, Potenzgesetze (Produkt gleicher Basen, Potenz einer Potenz), negative Exponenten als Kehrwert, Wurzel als Potenz mit Exponent ein Halb – selbst rechnen}
\begin{teile}
\teil $(5x)^3$ ohne Klammer? \leerfeld
\end{teile}
\end{aufgabe}
